\subsection{The exact class-three local order at two}

\begin{statement}{B.36}{The cubic determinant-extension class at two has exact order twenty-seven}{Proved}
Keep the $C_0$ tangent, primitive vector $w$, and actual
integral extension class $\eta$ of B.31 and B.35.
Then
\[
                              m_2=3.                         \tag{B.36.1}
\]
Consequently the actual zero-abelian class-two section
$a\kappa_3(2)w$, $a\in\mathbb Z_3$, lifts to intrinsic
class three if and only if $27\mid a$.
The same assertion holds after adding any
$z_0\in H^1(G_{\mathbb Q,\{3,\infty\}},\Lambda^2V)$.
In particular the coefficients $1,3,9$ do not lift,
whereas $27$ does.

This calculation uses the actual tangential translation
paths on the marked reduction, including their frames at
points of the modulus. Every possible integral Kummer
comparison is retained. The resulting finite affine sets
have nine elements for each of three comparisons; all
$27$ candidate classes have exact order $27$.
The proof computes an integral local torsion class and
does not determine the rational arithmetic Ross coefficient.
\end{statement}
\begin{proof}
\emph{The equation for an automorphism which moves the tangent.}
For every geometric automorphism $g$ of $(E_0,D)$ over
$\mathbb F_4$, choose a geometric path from the original
tangent to its image under $g$. Let $\widetilde g$ be the
resulting based automorphism. If its action fixes $w$, the
same commutator calculation as (B.35.7)--(B.35.8) gives
\[
 (g-I)[\eta]=\operatorname{ad}_{w*}(\alpha_g)
       \quad\hbox{in }H^1(\mathbb F_4,A(-1)).                \tag{B.36.2}
\]
Here $\alpha_g\in H^1(\mathbb F_4,V)$ is the abelian
class of that actual tangential path.
Indeed $\widetilde g\omega-\omega\in A$, and its
Frobenius difference is a coboundary; the inner
commutator contributes precisely $[\alpha_g,w]$.
Changing the path changes the representatives by these
coboundaries and does not change the displayed identity.
No finite-order assertion about the chosen based lift
is needed.
For the quaternion group this path class is zero;
for $\sigma$ it is $\kappa_3(\zeta)e_O$, as in B.35.
We compute it explicitly for the two translations by
$P_1=A_0$ and $P_2=A_2$.

\paragraph{The actual framed point group.}
Put $k=\mathbb F_4=\mathbb F_2[\zeta]$,
$E_0:y^2+y=x^3$, and let $D=E_0(k)=E_0[3]$.
The coordinate change and ordered cusps are those of B.35.
The inverse change is
\[
                  [X:Y:Z]=[y+\zeta:y+\zeta^2:x].
\]
Write $O$ for infinity.  For every $P\in D$ let $z_P$
be the point of the generalized Jacobian represented by
$\mathcal O(P-O)$ with the prescribed frames
$s_P^{-1}$ at $P$, $s_O$ at $O$, and one elsewhere;
put $z_O=0$.
The split torus is $(k^\times)^D/k^\times$, whose
coordinates, normalized at $O$, are written as
$W_0=\mathbb F_3^8$ using $\zeta$.
For $R\in D$ write $\epsilon_R(a)$ for the torus vector
having exponent $a$ at $R$ and zero elsewhere, modulo
the diagonal; in particular
$\epsilon_O(a)=(-a,\ldots,-a)$.
Every generalized-Jacobian point has a unique expression
$z_P+t$, with $P\in D$ and $t\in W_0$.
Indeed its compact image determines $P$, and the kernel
is the displayed split torus.  Surjectivity onto $E_0(k)$
is also explicit from the framed representatives $z_P$.

The point group is obtained without any choice of a
Tate division.  For $P,Q\ne O$ put $R=P+Q$.
If $R\ne O$, let
\[
 f_{P,Q}=\ell_{P,Q}/(x-x_R),
\]
where $\ell_{P,Q}$ is the secant or tangent line with
coefficient one on $y$.
If $R=O$, put $f_{P,Q}=x-x_P$; if either argument is
$O$, put $f_{P,Q}=1$.
In every case
\[
          \operatorname{div}(f_{P,Q})=P+Q-R-O.
\]
Multiplication by this function trivializes the
corresponding principal line bundle.  It sends its
prescribed frame at a cusp $T$ to the leading value
of $f_{P,Q}$ in $s_T$.  Therefore
\[
 z_P+z_Q-z_R=c(P,Q),\qquad
 c(P,Q)_T=
 \log_\zeta\frac{\operatorname{lead}_{s_T}f_{P,Q}}
                         {\operatorname{lead}_{s_O}f_{P,Q}}.
                                                               \tag{B.36.3}
\]
The use of frames is essential here: although $P$ and
$O$ belong to the modulus, $\mathcal O(P-O)$ is locally
free at both points. Its specified local generators
$s_P^{-1}$ and $s_O$ give nonzero elements of the
corresponding one-dimensional fibres. Thus the framed
bundle is well-defined, and a function with a zero or
pole is evaluated by its regularized leading coefficient.
This proves the sign of the toric correction.

Here are complete finite formulas for these leading
values.  At a finite cusp $T=(x_T,y_T)$ set
$\delta=x-x_T$.  The curve equation gives
\[
 y-y_T=x_T^2\delta+\delta^3+\delta^6+\cdots.
                                                               \tag{B.36.4}
\]
For a nonvertical line of slope $m$, its order and
leading value are read first from its constant term,
then from $x_T^2+m$, and finally from its cubic
coefficient one.  A vertical line has its ordinary
constant term or leading term $\delta$.
The actual parameters from the original plane model
are
\[
 s_{A_i}=(y+\zeta)/x,\quad
 s_{B_i}=(y+\zeta^2)/x,\quad
 s_{C_i}=x/(y+\zeta).
\]
Their derivatives relative to $dx$ are
\[
 k_T=\frac{ds_T}{dx}(T)=
 \begin{cases}
 x_T,&T=A_i\text{ or }B_i,\\
 (y_T+\zeta)^{-1},&T=C_1,C_2,\\
 1,&T=O.
 \end{cases}                                                   \tag{B.36.5}
\]
For the last equality, put $q=x/(y+\zeta)$ and
substitute $y=x/q+\zeta$ into the cubic equation.
It gives $x=q^{-2}+q+O(q^7)$, so $dx/dq$ has leading
coefficient one.  Also $y$ has leading term $q^{-3}$.
Thus all the functions in (B.36.3) have leading coefficient
one at $O$.
If their order at $T$ is $r$, multiply their leading
coefficient in $\delta$ by $k_T^{-r}$ to obtain their
leading coefficient in the prescribed parameter.

The program \path{c_char2_translation_paths.py} performs
all these operations in $k$ itself.
It checks the cocycle identity for $c$ on all $729$
triples and its symmetry on all $81$ pairs.
Consequently
\[
 (P,a)+(Q,b)=(P+Q,a+b+c(P,Q))                                \tag{B.36.6}
\]
is the complete framed group law, with $3^{10}$ elements.
For the compact basis $P_1=A_0$, $P_2=A_2$, it verifies
\[
                    3z_P=2U\,\operatorname{coord}(P)
                         \quad\text{in }W_0.                \tag{B.36.7}
\]
The two columns of $U$ are independent modulo three.
Hence $3J_{E,D}(k)$ has rank two over $\mathbb F_3$,
and (B.36.6) has type $C_9^2\oplus C_3^6$.

\paragraph{Automorphisms and their tangent paths.}
For a geometric automorphism $g$ in the subgroup
$E_0[3]\rtimes(Q_8\rtimes C_3)$ of order $216$, write
$g^*dx=c_gdx$.  The pushed local parameter at $gP$
is $s_P\circ g^{-1}$, whose derivative relative to
$s_{gP}$ is $k_P/(c_gk_{gP})$.
Thus the exact pushforward of the framed point is
\[
 \begin{split}
 g(z_P)={}&z_{gP}-z_{gO}
       +\epsilon_{gP}\left(\log_\zeta
                              \frac{c_gk_{gP}}{k_P}\right)\\
       &+\epsilon_{gO}\left(\log_\zeta
                              \frac1{c_gk_{gO}}\right).
 \end{split}                                                   \tag{B.36.8}
\]
On the torus $g$ permutes the nine coordinates before
renormalizing the $O$ coordinate.
Formula (B.36.8) is verified on all group laws for the
eight quaternion elements, the order-three generator,
and the two generating translations: $891$ checks.
It supplies genuine framed automorphism data, not
only cusp permutations.

For translation by $d$, $c_g=1$ and the image of the
original unit tangent has derivative $k_d$ at $d$.
Here is the tangential specialization with both endpoints
in the modulus. Over $k((\varepsilon))$ approach $d$ by
$P(\varepsilon)$ with $s_d(P)=a\varepsilon$, and approach
$O$ by $Q(\varepsilon)$ with $s_O(Q)=\varepsilon$.
The ordinary Abel divisor $P(\varepsilon)-Q(\varepsilon)$
is disjoint from $D$, so its canonical frames are one.
At $d$, that frame is $-a\varepsilon$ times the generator
$(s_d-a\varepsilon)^{-1}$. At $O$, it is
$(-\varepsilon)^{-1}$ times the generator
$s_O-\varepsilon$.
Use the compatible Frobenius-fixed roots
$\varepsilon^{1/3^n}$ of the tangential fibre to divide
the torus factors $\varepsilon$ at $d$ and
$\varepsilon^{-1}$ at $O$.
Reduction leaves the frames $a s_d^{-1}$ and $s_O$,
since minus one equals one in characteristic two.
The discarded formal principal units and the formal
reduction kernel are uniquely three-divisible, so their
Kummer classes vanish. The actual tangential class is
therefore $\kappa(z_d)+\kappa(a)e_d$, with the positive
sign. This calculation uses nonzero line-bundle frames
at the points of $D$, not meromorphic evaluation at a pole.
For the translated tangent $a=k_d$, its framed class is
\[
                         z_d+\epsilon_d(\log_\zeta k_d).
                                                               \tag{B.36.9}
\]
For an origin-fixing automorphism the compact divisor
is zero and the identical frame calculation gives the
toric class of its tangent multiplier.  In particular
it is zero for $Q_8$ and is
$\kappa_3(\zeta)e_O$ for $\sigma$, agreeing with
the independently proved tangent-rescaling identity.

Arithmetic Frobenius two is $(x,y)\mapsto(x^2,y^2+1)$.
The original plane parameters are defined over
$\mathbb F_2$, so it sends $z_P$ to $z_{F_2P}$ with
no additional frame multiplier; on the torus it squares
the values and permutes the cusps.
All $81$ corresponding group-law identities are checked.

\paragraph{The remaining three Kummer comparisons.}
Let $H=V/(F_4-I)V$ in the fixed integral lattice, with
\[
           F_4-I=\begin{pmatrix}-3I_2&0\\-2U&3I_8\end{pmatrix}.
                                                               \tag{B.36.10}
\]
A pair of integer vectors $(a,b)$ has unique normal form
\[
       (a\bmod3,\ b-2U\lfloor a/3\rfloor\bmod3)
                                                               \tag{B.36.11}
\]
in this finite group.
The ordinary Abel map induces the identity on compact
homology; the preceding local expansion induces the
positive identity on every peripheral generator.
Comparison of the two homology exact sequences thus
identifies $V$ with $T_3J_{E,D}$, including its integral
boundary and compact maps. An abelian tangential path
corresponds to precisely the framed Kummer class just
computed. The ordinary characteristic-zero identification
and its toric sequence are recorded in
\citet[\S1, immediately after (*) and (**)]{BSaidi2021};
their use at two is through the prime-to-two specialization
of \citet[XIII, \S\S2.9--2.12]{BSGA1} and the explicit
normalization just proved.
On the finite
field, Kummer gives the isomorphism
$J_{E,D}(k)\{3\}\simeq V/(F_4-I)V$: choose coherent
three-power divisions of a rational framed point and
take their Frobenius differences. Independence of those
divisions is exactly the coinvariant relation.
Here the entire point group is three-primary.
The canonical local Kummer isomorphism
$J_{E,D}(k)\xrightarrow{\sim}H$ sends a torus value
$\zeta$ to its corresponding peripheral coordinate.
For the compact sign, choose a coherent Tate sequence
$t_{i,r}\in E_0[3^r]$ with $t_{i,1}=P_i$ and
$3t_{i,r+1}=t_{i,r}$, representing $f_i$ modulo the
boundary. A $3^r$-division of $P_i$ is
$Q_r=t_{i,r+1}$. Since $F_4=[-2]$, its Kummer
Frobenius value is
$F_4Q_r-Q_r=-3Q_r=-t_{i,r}$.
Thus the compact Kummer class is $-f_i$, with its sign
fixed by the actual division equation.

Every comparison respecting these toric and compact
maps is therefore determined by
\[
                 z_{P_i}\longmapsto-f_i+b_i,
                         \qquad b_i\in\mathbb F_3^8.
                                                               \tag{B.36.12}
\]
Conversely (B.36.7) and (B.36.10) show that every such pair
defines a group isomorphism with the required graded
maps.  Equivariance under each actual automorphism
and Frobenius two imposes linear equations on the
sixteen entries of $b_1,b_2$.
The script constructs these equations from
(B.36.6)--(B.36.12), using the compatible quaternion,
CM, and Frobenius matrices of B.35.
It retains an integral $-1$ in every boundary column
which maps a cusp to $O$ before dividing the cross
block by three; replacing it early by the residue
$2$ would change the equation.

Exact elimination over $\mathbb F_3$ gives rank fifteen,
with particular solution
\[
 b^{(0)}=(1,2,2,2,1,2,1,1;\ 0,2,0,0,2,2,2,0)
\]
and kernel generator
\[
 b^{(1)}=(0,1,2,0,2,1,2,1;\ 1,2,0,2,1,0,2,1).
                                                               \tag{B.36.13}
\]
Thus the actual comparison is one of precisely
$b^{(0)}+a b^{(1)}$, $a\in\mathbb F_3$.
For each of them (B.36.9) gives an explicit actual-path
candidate in the fixed coefficient module.
The assertion is an exhaustive set of three possible
comparisons, not a claim that one arbitrarily selected
solution is the canonical Kummer map.
We retain all three in the following affine calculation.


\emph{The actual translation matrices preserve $w$ integrally.}
For $i=1,2$ let $P_i^{\rm bd}$ be the matrix of translation
by $P_i$ on the boundary lattice in its ordered eight
coordinates. A column whose image is $O$ is the vector
$(-1,\ldots,-1)^t$; every other column is the appropriate
unit vector. On compact homology a translation is the
identity. Commuting with the Frobenius-four matrix of
B.35 forces its degree-one matrix to be
\[
 T_i=\begin{pmatrix}
 I_2&0\\ (U-P_i^{\rm bd}U)/3&P_i^{\rm bd}
 \end{pmatrix},\qquad
 T_i\binom{3I_2}{U}=\binom{3I_2}{U}.                         \tag{B.36.14}
\]
The numerator is divisible by three because a translation
is the identity on the compact three-torsion quotient.
The displayed equality proves that both integral vectors
$u=3f_1+U_1$ and $v=3f_2+U_2$ are fixed. Thus
$T_iw=w$ exactly, as required for (B.36.2).
The scripts also check $T_i^3=I$, $T_1T_2=T_2T_1$, and
their commutation with Frobenius four.
Put $L_i=\operatorname{Lie}_3(T_i)$ and let
$S_A=\operatorname{Lie}_3(\sigma|V)$.
No action on the lower central series is inferred only
from a permutation: these are the functorial actions on
its actual integral degree-three graded piece.

\emph{The complete finite affine system.}
Let $D=4^{-1}\operatorname{Lie}_3(F_4|V)-I$ on the
rank-$330$ lattice $A(-1)$. Let $E$ be the integral
rank-$41$ basis of its quaternion-invariant direct
summand, and write $R_2$ for Frobenius two on this
summand. Use the Smith coordinates of B.35 for
\[
 C_2=\bigoplus_{j=1}^{13}\mathbb Z/3^{e_j},\qquad
 (e_j)=(1,1,1,1,1,2,2,3,3,3,3,4,4).                          \tag{B.36.15}
\]
If $Q_2$ is the stored row transport for $R_2-I$, let
$Q_2^{-1,+}$ denote its thirteen columns corresponding
to positive Smith exponents. Restriction to
$\mathbb F_4$ is represented by
\[
                   B=E(I+R_2)Q_2^{-1,+}.                     \tag{B.36.16}
\]
Restriction is injective and preserves orders, since
corestriction composed with restriction is multiplication
by the unit two. Averaging over $Q_8$ commutes with Frobenius
two, because that Frobenius normalizes $Q_8$.
Hence the actual class occurs among these thirteen
coordinates.

For each comparison $a\in\mathbb F_3$ of (B.36.13),
let $\alpha_i^{(a)}$ be (B.36.9) evaluated with that
comparison. The complete set tested here is
\[
 \begin{aligned}
 (S_A-I)Bx&=\operatorname{ad}_w(e_O),\\
 (L_i-I)Bx&=\operatorname{ad}_w(\alpha_i^{(a)}),
                &&i=1,2,
 \end{aligned}
 \qquad\pmod{DA}.                                          \tag{B.36.17}
\]
All three identities are necessary for the actual
arithmetic extension, by (B.36.2). The first is B.35's
independent tangent-rescaling identity, with
$\kappa_3(\zeta)(F_4)=1$ modulo three.
The targets are taken in the full rank-$330$ cokernel:
a translation need not preserve the quaternion-invariant
subspace. Thus no target component is discarded.

Here is an explicit small certificate of the complete
solution set. In coordinates (B.36.15), separated after
positions five and eleven, put
\[
 \begin{aligned}
 b={}&(0,0,1,2,0;\ 2,5,9,6,9,24;\ 3,36),\\
 h_3={}&(0,0,2,1,1;\ 3,0,0,18,9,0;\ 0,0),\\
 h_9={}&(0,0,0,0,0;\ 6,6,0,18,0,18;\ 9,27).
 \end{aligned}                                               \tag{B.36.18}
\]
The subscripts are their exact orders. For comparison
$a$ the entire solution set of (B.36.17) is
\[
 \left\{b+s h_3+t h_9:\ s\in\{0,1,2\},\
               t\in\{0,\ldots,8\},\ t\equiv a\pmod3\right\}.
                                                               \tag{B.36.19}
\]
There are nine distinct elements for each $a$.
Its homogeneous difference group is generated by $h_3$
and $3h_9$, and is $(\mathbb Z/3)^2$.
Every one of the $27$ displayed candidates satisfies
\[
                         27x=0,\qquad x_{12}\equiv3\pmod9.
                                                               \tag{B.36.20}
\]
The twelfth coordinate is in $\mathbb Z/81$, so these
conditions prove that its exact order is $27$.
They do not require selection of a particular canonical
Kummer comparison or a particular extension representative.

For reproducibility, we describe the finite elimination
which proves completeness of (B.36.19), in addition to
checking all its members. In the integral Hall basis
$[[e_i,e_j],e_k]$ with $i>j$ and $k\ge j$, Smith reduction
of the full operator $D$ gives the valuations
\[
                      0^{132},\ 1^{70},\ 2^{24},\ 3^{76},\ 4^{28}.
                                                               \tag{B.36.21}
\]
Write $Q$ for its recorded invertible row transport.
For $g=S_A,L_1,L_2$, apply $Q$ to $(g-I)B$ and to its
right side in (B.36.17), discard only the $132$ zero
Smith factors, and reduce each remaining row modulo
its corresponding power of three. This gives $594$
explicit scalar congruences in thirteen variables.
The code checks that multiplying column $j$ by $3^{e_j}$
annihilates each row, so the equations are well-defined
on $C_2$.

The rowwise solver maintains the exact affine lattice
$x=x_0+H y$. For a new row $c x=d\pmod{3^r}$,
substitution gives $(cH)y=d-cx_0$. Let $3^v$ be the
common three-power divisor of the entries of $cH$,
capped at $3^r$. If the residual is not divisible by
$3^v$, the system is inconsistent. Otherwise divide
by $3^v$, choose a unit coefficient, and solve for that
coordinate modulo $3^{r-v}$. Substitution into the other
coordinates and multiplication of its free column by
$3^{r-v}$ gives exactly all solutions of the enlarged
system. Induction proves completeness of this algorithm.
All computations can be reduced modulo $81$, because
both the domain and target groups are annihilated by
$81$.
The resulting homogeneous lattice has order $9$ and
exponent $3$, independently agreeing with the dual
coinvariant Smith calculation in
\path{b_char2_translation_linear.py}.
For each comparison the affine solver gives nine
solutions, verifies all $27$ against every row, and
checks equality with (B.36.19).
It independently imposes $3^k x=0$ for $k=1,2,3,4$;
the four feasibility answers are false, false, true,
true for every comparison.
An independent verifier, \path{c_char2_translation_verify.py},
uses ternary digit lifting instead of affine-lattice
elimination. Its successive candidate counts are
$243,2187,19683,59049$. It enumerates all $3^{10}$
solutions of the original quaternion, sigma, and
Frobenius-two conditions, then tests every candidate
against both translation equations for each comparison.
It finds exactly the same nine solutions in each case,
with all $27$ orders equal to $27$.

The matrix entries involve only the Hensel root
$s^2=-2$, $s\equiv1\pmod3$, from B.35.
The program computes this root modulo $3^{11}$ before
the integral divisions by three, obtaining all action
matrices modulo $3^{10}$. Every Smith pivot just listed
is strictly below this precision, so the cokernel and
all its affine equations modulo $81$ are exact.
The executable certificates are
\path{c_char2_translation_paths.py},
\path{b_char2_translation_linear.py}, and
\path{b_char2_translation_affine.py}, together with their
recorded matrices, outputs, and the explicit $27$-entry
JSON list. The first reconstructs the framed group and
all three comparisons; the latter two retain the actual
integral action and all path terms.

The actual local extension satisfies (B.36.17) for one
of the three exhaustive comparisons, hence has order
$27$ by (B.36.19)--(B.36.20). Restriction and quaternion
averaging preserve its order, proving (B.36.1).
Finally B.31's exact global local-order and decay law
applied to support at two gives $27\mid a$ as both a
necessary and a sufficient condition for the stated
class-three lift. A single support prime satisfies the
decay condition automatically, and its unramified
summand $z_0$ has zero obstruction as proved there.
\end{proof}
